\chapter{Appearance Is Remembered}

\section{The Second Glance}

Chapter 8 left the courtyard unresolved on purpose. Its content was live, uncertain, entropy high, and nothing in the argument so far said what happens the moment someone actually walks around the corner and looks. That moment is where this chapter begins.

Suppose the observer looks, and the field, through whatever process eventually generates a concrete perception from an uncertain region, produces a courtyard with a stone bench near the far wall and a scattering of gravel underfoot. This act of generation is not a violation of anything argued so far. It is, in fact, required: the courtyard had to become something specific eventually, and nothing about Chapter 8's entropy constraint forbids resolution at the moment of observation, only unearned resolution beforehand. A bench and some gravel, produced now, for the first time, are a legitimate answer to a genuinely open question.

Now suppose the observer walks away, comes back a minute later, and looks again. What must be true of the second glance for this to be the same courtyard rather than merely a second, unrelated act of invention? The bench must still be there, in the same place. The gravel must still be gravel, not suddenly flagstone. Nothing about the second glance is permitted to reopen what the first glance already closed. This is a stronger requirement than it might first appear, because a generative system built the ordinary way, sampling plausibly from a learned distribution each time it is asked, will happily produce a coherent, high-quality courtyard on the second glance that has nothing to do with the courtyard the first glance produced. Both would be individually convincing. Neither would be faithful to the other. This is precisely the failure this book set out from in Chapter 6, now returned to at the one point in the field's dynamics where it actually occurs: the instant an uncertain region is first observed and must decide what it is going to remain.

\section{The Two Kinds of Generation}

It is worth separating two things that look identical from the outside but occupy entirely different roles in this framework, because collapsing them is exactly the mistake that produces the wall problem in the first place.

The first is generation as invention. This happens exactly once for any given region of high entropy, at the moment an observation first presses on it. Before that moment, nothing in the field specifies a bench rather than a fountain rather than an empty stretch of paving; the courtyard's entropy was genuinely high, and something has to break the symmetry. Invention of this kind is not merely permitted, it is unavoidable, and there is nothing dishonest about it. A system that refused to invent anything until it had somehow already been told the answer would never be able to observe anything for the first time at all.

The second is generation as regeneration, and this is the pathology. It occurs when a system treats every observation as an occasion for a fresh act of invention, whether or not the region in question has already been resolved. Regeneration looks, locally, exactly like invention: the same decoder, the same renderer, the same plausible output. The difference is entirely a matter of what should have happened instead. Invention answers a question that was genuinely still open. Regeneration re-answers a question that had already been closed, discarding the previous answer without acknowledging that it ever existed.

The entire discipline this chapter is about to formalize exists to keep the first kind of generation legitimate while making the second kind of generation structurally impossible, not merely discouraged. This cannot be achieved by training a model to prefer consistency, because a model that merely prefers consistency will still, on a long enough encounter, invent a courtyard incompatible with its own earlier output and call it plausible. It has to be achieved by giving the field somewhere to keep what invention has already produced, so that the second glance is reading a memory rather than performing a second act of invention that happens to resemble the first.

\section{The Invention-to-Memory Principle}

That somewhere is z, the appearance component introduced in Chapter 7 and left, until now, largely undeveloped. Its update rule is the mechanism this chapter has been building toward:

\[ z_{t+1}(x) = z_t(x) + \lambda \cdot \hat{z}(x). \]

Here ẑ(x) is a candidate: a proposed appearance for the point x, produced by whatever process generates candidates, a decoder, a renderer, a language model completing a description, when an observation presses on that point. It is not yet committed. The coefficient λ determines how much of that candidate is actually allowed to update the field's held value at x, and its behavior over time is what separates invention from regeneration. When x is genuinely unresolved, when entropy there is still high and nothing has previously been committed, λ can be large: the candidate is mostly accepted, because there is little existing content for it to conflict with, and this is invention doing its legitimate work. Once x has been observed and \(z_t(x)\) already holds a specific, settled value, the same update rule behaves very differently. A new candidate ẑ(x) arising from a second observation is not free to overwrite what is already there; the update is constrained to move \(z_t(x)\) only as far as continued fidelity to the existing value permits, and in the limiting case, once a region is fully settled, λ's effective contribution collapses toward nothing. The bench does not become a fountain on the second glance, not because the renderer forgot how to imagine a fountain, but because the update rule no longer gives an unconstrained candidate the power to overwrite an already-committed one.

This is the invention-to-memory principle stated in full: a region may invent its appearance once, while its entropy is genuinely open, and having invented it, the field is obligated to hold that appearance stable against future candidates rather than regenerating it. The principle does not forbid change entirely. A wall can be damaged, a bench can be moved, and Part III's treatment of proposals and constraint projection will have a great deal to say about how legitimate changes to an already-settled region are distinguished from illegitimate ones. What the invention-to-memory principle forbids specifically is change with no cause: appearance drifting simply because it was asked about again, with nothing in the intervening history to justify a different answer.

\section{Appearance and Entropy, Together}

It is worth making explicit a connection that the last two chapters have been circling without quite stating outright. Chapter 8's entropy constraint said that S may not fall at a point unless that point has been observed. This chapter's invention-to-memory principle says what happens at exactly the moment that license is used: entropy falls, and z is what it falls into. The two are not separate mechanisms that happen to cooperate. They are two descriptions of a single event. To observe an uncertain region is to reduce its entropy and commit its appearance in the same act, and a field that did only the first without the second would have no way of remembering what its own reduced uncertainty had resolved into; a field that did only the second without the first would be committing appearances without any record of how settled or provisional they actually were. The five-component decomposition from Chapter 7 earns its keep exactly here, at the seam between S and z, where the honest bookkeeping of what remains uncertain and the honest bookkeeping of what has already been decided turn out to be two sides of a single operation.

\section{Rendering as a Concrete Observation Operator}

Chapter 6 introduced \(O_\theta\) only as an abstract projection from field to observation, and Chapter 7 described it only as a family of operators, each recovering some partial slice of the field's structure. It is now possible to say something concrete about at least one member of that family, because appearance is the component most observation operators are actually built to read.

A renderer, in the ordinary sense, computes an approximation to \(O_\theta\) specialized to z: given the field's current appearance at and around a point, it produces an image, a sound, a description, whatever the modality requires. This is the forward direction, field to observation, and it is the direction Chapter 6 emphasized. But the invention-to-memory principle depends equally on the reverse direction. The candidate ẑ(x) in the update rule has to come from somewhere, and it comes from an operator that runs, loosely speaking, backward: given an observation, or given the field's surrounding context at a point whose appearance is not yet settled, produce a proposal for what that point's appearance should be. This is inverse rendering in the broad sense this book will use the term, and diffusion decoders, language models completing a scene description, and physically grounded renderers run in reverse are all instances of it. The forward operator reads what the field already holds. The inverse operator proposes what the field might come to hold. Between them, they are the two halves of the observation-and-invention cycle this chapter has been developing, and Chapter 27 will return to both in their fully engineered form, as concrete software components rather than the conceptual operators sketched here.

It is worth being honest about what this section has and has not established. It has said that rendering and inverse rendering are the concrete faces of \(O_\theta\) where appearance is concerned, and it has said that the candidate driving the invention-to-memory update comes from the inverse direction. It has not said how either operator is actually built, what architecture a renderer or an inverse renderer should use, or how the two are kept consistent with one another in practice. Those are engineering questions, and this book's four-pass structure defers them deliberately, the same way Chapter 6 deferred the mathematics of constraint projection to Part IV. What matters here is only that \(O_\theta\), introduced three chapters ago as an unspecified projection, has now acquired its first genuinely concrete instance.

\section{Looking Ahead}

Appearance answers what a region currently looks like, and this chapter has explained why that answer, once given, must be kept rather than regenerated. But looking like something and being usable are not the same question. A bench that has settled into the field with a fixed appearance affords sitting to a human passerby and affords nothing in particular to whatever else might encounter it. The field's appearance is now well defined; what a given observer is entitled to do with it is not yet addressed at all. That is the question Chapter 10 takes up, and it is the question that will finally put to use the capability-relative structure this book set aside back when the six-component decomposition was refined into five.
